% ================================================================================
% Section 3: The Modulo Trick - Intuitive Introduction
% ================================================================================
\section{Understanding the Modulo Trick}

\begin{frame}{What is the Modulo Trick? An Intuitive Approach}
    \begin{itemize}
        \item Imagine you're trying to guess a number, but you can only ask questions about its remainder when divided by certain numbers.
        \item \textbf{Example:} I'm thinking of a number $x$.
        \begin{itemize}
            \item If I tell you $x \equiv 1 \pmod 2$ (odd), what do you know?
            \item If I tell you $x \equiv 2 \pmod 3$, what do you know?
            \item Combining both, you can narrow down possibilities!
        \end{itemize}
        \item This is exactly what the modulo trick does - it uses remainders to restrict possible values of exponents.
    \end{itemize}
\end{frame}

\begin{frame}{The Basic Idea of the Modulo Trick}
    \begin{block}{Core Concept}
        Instead of solving $1+p^a = q^b + r^c$ directly, we ask:
        "What must be true about $a,b,c$ modulo different numbers?"
    \end{block}

    \vspace{0.3cm}
    \textbf{Why does this help?}
    \begin{itemize}
        \item Exponential expressions have predictable patterns modulo $m$
        \item By choosing the right moduli, we can force contradictions unless exponents are small
        \item Once exponents are bounded, we can check all remaining possibilities by brute force
    \end{itemize}
\end{frame}
